% !PRINTED 606
\weberpage{606}
%<<CONT
und wenn man bedenkt, dass die Grundzahl $\Delta$ im Falle a) gleich $4d$, im Falle b) gleich $d$ ist, so erhalten wir, wenn wir
\weq{5}{x = 2a_{1,2} \text{ oder } {}= 2a_{1,2}+1}
setzen, folgendes Resultat:

1. \gesperrt{Die nothwendige und hinreichende Bedingung daf\"ur, dass $p$ zwei Primfactoren enth\"alt, ist die, dass die Congruenz}
\weq{6}{x^2 \equiv \Delta\ (\mathrm{mod}\ 4p)}
\gesperrt{durch eine ganze rationale Zahl $x$ befriedigt werden kann, oder dass die Grundzahl $\Delta$ quadratischer Rest von $4p$ ist.}

Bezeichnen wir mit $\pi, \pi'$ die beiden conjugirten Formen
\weq{7}{\pi = pu + (a_{1,2}+\Theta)v,\quad \pi' = pu + (a_{1,2}+\Theta')v,}
so ergiebt sich durch Multiplication
\weq{8}{\pi\pi' = p\left(pu^2 + xuv + \frac{x^2-\Delta}{4p}v^2\right).}

Hierin ist nun
\[
pu^2 + xuv + \frac{x^2-\Delta}{4p}v^2
\]
eine Einheit, weil, wenn $p$ in $\Delta$ nicht aufgeht, $x$ durch $p$ nicht theilbar ist, und wenn $p$ in $\Delta$ aufgeht, zwar $x$, aber nicht $(x^2-\Delta:4p)$ durch $p$ theilbar ist. Denn ist $p$ ungerade, so kann $p^2$ nicht in $\Delta$ aufgehen, und ist $p=2$ und $\Delta$ durch 2 theilbar, so ist $\Delta = 4d$, $d \equiv 2, 3\ (\mathrm{mod}\ 4)$. W\"are aber
\[
x^2 - \Delta \equiv 0\ (\mathrm{mod}\ 16),
\]
so w\"urde folgen
\[
d \equiv \left(\frac{x}{2}\right)^2 (\mathrm{mod}\ 4),
\]
w\"ahrend doch $\left(\dfrac{x}{2}\right)^2$ als Quadrat nicht mit 2 oder 3 congruent sein kann. Demnach ist $p$ in die beiden Primfactoren $\pi$ und $\pi'$ zerlegt.

Die beiden Factoren $\pi$ und $\pi'$ von $p$ k\"onnen auch mit einander associirt sein. Dies tritt nur dann ein, wenn $\pi'$ und folglich auch $\pi - \pi'$ durch $\pi$ theilbar ist. Nun ist aber nach (7)
\[
\pi - \pi' = (\Theta - \Theta')v,
\]
und folglich muss
\[
(\Theta - \Theta')^2 = \Delta
\]
%CONT>>
